\documentclass[12pt]{article}
\usepackage[utf8]{inputenc}
\usepackage[T1]{fontenc}
\usepackage{hyperref}
\title{Tewa Language Resources for K-12 Curriculum Development}
\author{Compiled by Claude Sonnet 5 (claude-sonnet-5) for Andr\'as Kornai}
\date{}
\begin{document}
\maketitle

\section{Introduction}
Tewa (ISO 639-3 tew; Glottocode tewa1260) is a Kiowa-Tanoan language of the Pueblo communities of northern New Mexico, with about 1,500 speakers. Searches found only a small amount of publicly accessible material. Orthographies differ from pueblo to pueblo, which complicates shared materials. The Tewa pueblos have different views on writing the language, and some elders prefer oral transmission only. This report lists only what appeared in search results.

\section{Primary Educational Resources}
No publicly available K-12 textbook or primer was found. The Nambe Pueblo Tewa Language Program is listed as offering classes for all ages: \url{https://sharenm.org/organization/print/nambe-pueblo/nambe-pueblo-tewa-language-program}.
A Worcester Polytechnic Institute project report describes a cultural-education platform, with videos narrated in Tewa, that supplements an existing Tewa curriculum at a northern New Mexico pueblo: \url{https://digitalcommons.wpi.edu/iqp-all/3299}.

\section{Linguistic and Reference Materials}
The \emph{San Juan Pueblo Tewa Dictionary}, compiled by Esther Martinez with the San Juan Pueblo Bilingual Program, is the main reference. Sources disagree on the date (1982 or 1992). The UBC RelLex entry notes that it is print-only, without audio or example phrases: \url{https://wiki.ubc.ca/Documentation:RelLex/San_Juan_Pueblo_T\%C3\%A9wa_Dictionary}.

\section{Religious and Community Resources}
Biographical information on Esther Martinez, the Tewa speaker and language advocate: \url{https://en.wikipedia.org/wiki/Esther_Martinez}.
The San Ildefonso Pueblo law library has a council code entry that surfaced in the search. Its relevance to language policy was not verified: \url{https://law.sanipueblo.org/us/nsn/san-ildefonso/council/code/6.5.2.010}.

\section{Digital and Audio Resources}
No audio archive or app was found in these searches. The WPI platform above includes Tewa-narrated video, but it may not be publicly accessible.

\section{Curriculum Development Context}
Little published material is available. Any curriculum would have to be built with the individual pueblo language programs, such as Nambe and San Juan, and would need to settle which orthography to use. The programs should be contacted directly.
\end{document}